% !TEX program = xelatex
% Overleaf: Menu -> Compiler -> XeLaTeX
\documentclass[11pt]{article}

\usepackage[review]{acl}

\usepackage{latexsym}
\usepackage{microtype}
\usepackage{inconsolata}
\usepackage{graphicx}
\usepackage{booktabs}
\usepackage{multirow}
\usepackage{amsmath}

\usepackage{fontspec}
\usepackage{xeCJK}
\setmainfont{TeX Gyre Termes}
\setCJKmainfont{FandolSong-Regular}[
  BoldFont = FandolSong-Bold,
  ItalicFont = FandolSong-Regular,
  Scale = 0.95,
]

\title{The Illusion of Parity: Human Cultural Adequacy Reveals Shared Whitewashing in Dialectal Translation}

\author{Anonymous ACL submission}

\begin{document}
\maketitle

\begin{abstract}
Evaluating whether NLG outputs preserve \emph{variation}---register, pragmatics, and locally grounded meaning---is central to personalised generation, yet standard reference-based scores may not reflect such loss on dialectal input.
We investigate whether automatic metrics capture culturally grounded meaning in Sichuanese$\rightarrow$English translation on a 200-sentence keyword-stratified probe set from WenetSpeech-Chuan.
Two contemporary LLMs---GLM-4-Plus and DeepSeek-V4-Flash\footnote{We invoke the public DeepSeek Chat API (\texttt{deepseek-chat}); our provider dashboard reports usage under DeepSeek-V4-Flash.}---yield near-identical system-level COMET scores (70.76 vs.\ 70.87; paired bootstrap $p=0.836$), suggesting corpus-level parity.
Human \emph{cultural adequacy} judgments paint a different picture: only 14\% of sentences are adequate from \emph{both} models, whereas 28\% fail from \emph{both}---a pattern we term \emph{consensus whitewashing}, defined as the shared loss of culturally marked pragmatic meaning across systems.
McNemar's test on paired pass/fail labels favors DeepSeek-V4-Flash over GLM-4-Plus (34.0\% vs.\ 52.0\% cultural pass; $p=0.001$).
Automatic metrics partially separate successful from failed cases (${\approx}12$ COMET points between \emph{tie-good} and \emph{tie-bad}), yet failed outputs still score ${\approx}67$ COMET; sentence-level COMET correlates only weakly with human adequacy (Spearman $\rho=0.29$), and agrees with human winners on just 46.6\% of contested sentences.
We release probe-construction scripts and a lightweight $2{\times}2$ annotation protocol ($\kappa=0.71$) for variation-rich NLG evaluation.
\end{abstract}

\section{Introduction}

Personalised and dialectal NLG must preserve speaker-specific register and cultural pragmatics, not only fluent standard-language output \cite{hendy-etal-2023-mt}.
Regional varieties therefore stress \textbf{NLG evaluation} itself: corpus-level, reference-based scores designed for standard machine translation (MT) may rank systems similarly even when both outputs normalize culturally marked meaning---flattening emphatic self-reference, insult register, or domain-specific lexicon toward generic equivalents.
Dialectal translation differs from standard MT because success requires transmitting implicit socio-cultural context, discourse markers, and register that are often absent from high-resource benchmarks.

We describe this systematic loss as \textbf{cultural whitewashing}: an \emph{analytical} concept for the erosion of culturally grounded pragmatic meaning toward bland, globally standardized paraphrases---not a moral judgment about model intent.
Evaluating whitewashing requires moving beyond aggregate string or embedding similarity.

On our Sichuanese probe set, GLM-4-Plus and DeepSeek-V4-Flash achieve statistically indistinguishable BLEU and COMET (Table~\ref{tab:system}), creating an \textbf{illusion of parity}.
We address three research questions:
\textbf{RQ1:} Do aggregate automatic scores mask disparities in cultural preservation?
\textbf{RQ2:} Do keyword-stratified scores reveal pragmatic valleys invisible at system level?
\textbf{RQ3:} When models fail culturally, is failure shared across systems and under-penalised by reference-based metrics?

We combine (1)~keyword-enriched sampling over WenetSpeech-Chuan transcriptions;
(2)~system-level and keyword-stratified BLEU, chrF, and COMET with bootstrap tests;
and (3)~a blind $2{\times}2$ human \emph{cultural adequacy} protocol that labels each sentence by whether \emph{each} model transmits source cultural/pragmatic meaning---not pairwise fluency preference.

Our contributions are:
(1) a reproducible keyword-stratified Sichuanese probe from WSC-Train;
(2) a lightweight human protocol with substantial agreement ($\kappa=0.71$);
(3) quantitative evidence of metric--human divergence and consensus whitewashing; and
(4) a qualitative typology with implications for culturally informed NLG evaluation.

\section{Related Work}

\subsection{Dialectal Translation and Resources}
Translating dialectal speech presents challenges distinct from standard MT: dialects encode implicit socio-cultural context, discourse markers, and specialized lexicon.
WenetSpeech-Chuan (WSC-Train) provides large-scale Sichuanese transcriptions from in-the-wild speech \cite{wenetspeech-chuan}.
Models trained predominantly on standard Mandarin or English often normalize regional forms toward mainstream lexicon and syntax---precisely the variation-sensitive behavior NLG evaluation should surface.

\subsection{Automatic MT Evaluation}
LLMs translate competitively on high-resource pairs \cite{hendy-etal-2023-mt}, typically assessed with BLEU \cite{post-2018-sacrebleu}, chrF \cite{popovic-2015-chrf}, and COMET \cite{rei-etal-2022-comet}.
These metrics assume human references encode semantic targets worth preserving.
When references are culturally marked---as in our human-written Sichuanese--English gold---metrics may still reward generic fluent alternatives that align partially with reference wording while missing source pragmatics \cite{freitag-etal-2021-bleu}.

\subsection{Human Evaluation Beyond Fluency}
Expert and error-span annotations reveal systematic gaps between metric rankings and human judgments \cite{freitag-etal-2021-bleu}.
Our work extends this line to \emph{cultural adequacy} on dialectal input, using a four-way consensus scheme designed to detect \emph{shared} failure across models rather than simple pairwise preferences.

\section{Methodology}

\subsection{Evaluation Framework}
Our framework stratifies evaluation by dialect keyword type and assessment modality: a keyword-driven probe ensures dense dialectal phenomena, automatic metrics supply corpus-level scores, and human cultural adequacy tests whether each model transmits source pragmatics.
This design targets \emph{personalisation and variation} in NLG evaluation rather than a single leaderboard score.

\subsection{Probe Set Construction}

\paragraph{Source corpus.}
We download \texttt{wsc\_metadata.jsonl} from HuggingFace \texttt{ASLP-lab/WSC-Train} \cite{wenetspeech-chuan} and retain utterances with ASR confidence $\geq 0.9$, yielding 2{,}598{,}002 sentences.

\paragraph{Keyword stratification and sampling rationale.}
We construct a \textbf{probe set} ($n=200$, seed 42), not a random corpus sample.
Uniform sampling would dilute culture-bearing terms: high-frequency particles such as 啥子 dominate unconstrained data while rare insult/register items (e.g., 瓜娃子) seldom appear.
We therefore define 45 keyword patterns in three strata: (i) \emph{content-bearing} terms (老子, 安逸, 瓜娃子, 咋个); (ii) \emph{modal particles} (嘛, 哈, 哦); and (iii) \emph{long-tail cultural tokens} (e.g., mahjong 九通, 扯把子).
High false-positive single-character patterns use context constraints; filler tokens and ASCII fragments are removed.
Sentences must contain at least one content keyword and have length $\leq 80$ characters.
Weighted sampling uses content weight 1.0, particle weight 0.15, downweights 啥子, and upweights rare register terms---balancing diagnostic coverage with tractable human annotation.
The probe \textbf{stress-tests} variation-sensitive failure modes; prevalence estimates for WSC-Train as a whole are out of scope.

\paragraph{Reference translations.}
Two bilingual annotators independently authored culturally faithful English references for all 200 sentences.
References are used only for automatic metrics; models never see them at inference.

\subsection{Translation Systems}
\label{sec:systems}
Both models receive an identical naive zero-shot prompt:
\textit{``You are an expert translator.''} (system) and
\textit{``请将以下四川话翻译为地道的英文：\{text\}''} (user).
We evaluate GLM-4-Plus (\texttt{glm-4-plus}) and DeepSeek-V4-Flash (\texttt{deepseek-chat}) via public APIs.
We compare exactly these two systems; results do not generalize to broader open- vs.\ closed-source categories.

\subsection{Automatic Metrics}
We compute corpus-level BLEU and chrF (sacreBLEU) and sentence-level COMET (\texttt{Unbabel/wmt22-comet-da}, $\times 100$).
System-level differences use paired bootstrap resampling ($B=2000$; \citealp{koehn-2004-statistical}).
We report keyword-stratified chrF/COMET (minimum $n=5$ in main tables) and flag refusal outputs (GLM-4-Plus: 1/200; DeepSeek-V4-Flash: 0/200).

\subsection{Human Cultural Adequacy Protocol}

\paragraph{Task and criteria.}
Annotators view the Sichuanese source and two blinded English translations.
For each sentence they assign one consensus label based on whether each translation \emph{successfully transmits the cultural/pragmatic meaning} of the source---not general fluency alone (Table~\ref{tab:labels}).

\begin{table}[t]
  \centering
  \small
  \begin{tabular}{@{}cllp{4.2cm}@{}}
    \toprule
    \textbf{Label} & \textbf{GLM} & \textbf{DS} & \textbf{Meaning} \\
    \midrule
    \emph{tie-good} & \checkmark & \checkmark & Both culturally adequate \\
    \emph{g} & \checkmark & $\times$ & Only GLM-4-Plus adequate \\
    \emph{d} & $\times$ & \checkmark & Only DeepSeek-V4-Flash adequate \\
    \emph{tie-bad} & $\times$ & $\times$ & Consensus whitewashing \\
    \bottomrule
  \end{tabular}
  \caption{Human cultural adequacy labels ($n=200$). DS = DeepSeek-V4-Flash.}
  \label{tab:labels}
\end{table}

\paragraph{Workflow and agreement.}
Two bilingual annotators independently labeled all 200 sentences, then discussed disagreements to reach consensus gold labels stored in \texttt{preference}.
Pre-consensus Cohen's $\kappa=0.71$ indicates substantial agreement \cite{cohen1960kappa}.
Guidelines operationalise: 安逸 as Sichuanese pleasure/satisfaction (not merely ``comfortable''); 老子 as emphatic self-reference (not ``father'' or ``Laozi''); 瓜娃子 as insult register (not neutral ``silly person'').

\section{Results}

\subsection{RQ1: System-Level Parity and Human Pass Rates}

Table~\ref{tab:system} reports system-level scores.
DeepSeek-V4-Flash leads on all three automatic metrics, but only chrF differs significantly ($+1.98$, 95\% CI $[0.62, 3.40]$, $p=0.006$).
COMET ($+0.11$, $p=0.836$) and BLEU ($+0.20$, $p=0.861$) show no significant difference---the illusion of parity at corpus level.
Sentence-level oracle COMET (selecting the better model per sentence) reaches 73.35, indicating headroom neither model captures alone, without locating \emph{where} cultural meaning is lost.

Human cultural pass rates diverge: 34.0\% (GLM-4-Plus) vs.\ 52.0\% (DeepSeek-V4-Flash).
On paired binary judgments, 40 sentences are adequate only under GLM-4-Plus (\emph{g}) vs.\ 76 only under DeepSeek-V4-Flash (\emph{d}); McNemar's exact test on discordant pairs yields $p=0.001$.

\begin{table}[t]
  \centering
  \small
  \begin{tabular}{@{}lrrrr@{}}
    \toprule
    \textbf{Model} & \textbf{BLEU} & \textbf{chrF} & \textbf{COMET} & \textbf{Cult.\ pass} \\
    \midrule
    GLM-4-Plus & 16.12 & 39.75 & 70.76 & 34.0\% \\
    DeepSeek-V4-Flash & 16.32 & 41.73 & 70.87 & 52.0\% \\
    \midrule
    \multicolumn{5}{@{}l@{}}{\footnotesize Bootstrap $p$: COMET 0.836; BLEU 0.861; chrF 0.006. McNemar $p=0.001$.} \\
    \bottomrule
  \end{tabular}
  \caption{System-level automatic scores and human cultural pass rates ($n=200$).}
  \label{tab:system}
\end{table}

\subsection{RQ2: Keyword-Stratified Pragmatic Valleys}

Aggregate parity masks local deficits on content-bearing terms (Table~\ref{tab:keywords}).
老子 (emphatic self-reference) averages COMET ${\approx}66.7$ for both models---${\approx}4$ points below the corpus mean---with a 50.0\% \emph{tie-bad} rate (9/18).
安逸 (Sichuanese pleasure) averages COMET 68.3--69.1 with chrF as low as 34.0 (GLM).
High-frequency 啥子 ($n=44$) yields near-parity COMET (${\approx}71$), showing stratification separates content words from discourse markers.

\begin{table}[t]
  \centering
  \small
  \begin{tabular}{@{}lrrrrr@{}}
    \toprule
    \textbf{Keyword} & \textbf{$n$} & \textbf{GLM-C} & \textbf{DS-C} & \textbf{GLM-F} & \textbf{DS-F} \\
    \midrule
    老子 & 18 & 66.78 & 66.72 & 34.36 & 37.53 \\
    安逸 & 42 & 68.34 & 69.08 & 34.02 & 36.59 \\
    瓜娃子 & 5 & 76.94 & 69.82 & 46.14 & 41.89 \\
    咋个 & 11 & 71.98 & 75.26 & 46.16 & 51.49 \\
    啥子 & 44 & 70.71 & 71.40 & 40.29 & 41.70 \\
    \bottomrule
  \end{tabular}
  \caption{Keyword-stratified scores ($n\geq 5$). C = COMET ($\times 100$); F = chrF; DS = DeepSeek-V4-Flash.}
  \label{tab:keywords}
\end{table}

\subsection{RQ3: Consensus Whitewashing and Metric Divergence}

Table~\ref{tab:gold} summarizes gold labels: 28 \emph{tie-good} (14.0\%), 40 \emph{g} (20.0\%), 76 \emph{d} (38.0\%), and 56 \emph{tie-bad} (28.0\%).
The modal class is \emph{d}, yet over one quarter of sentences exhibit \emph{consensus} failure---whitewashing is often shared, not model-specific.

\begin{table}[t]
  \centering
  \small
  \begin{tabular}{@{}lrr@{}}
    \toprule
    \textbf{Label} & \textbf{Count} & \textbf{Percent} \\
    \midrule
    \emph{tie-good} (both adequate) & 28 & 14.0\% \\
    \emph{g} (GLM only) & 40 & 20.0\% \\
    \emph{d} (DeepSeek only) & 76 & 38.0\% \\
    \emph{tie-bad} (consensus failure) & 56 & 28.0\% \\
    \bottomrule
  \end{tabular}
  \caption{Human cultural adequacy gold distribution ($n=200$).}
  \label{tab:gold}
\end{table}

Table~\ref{tab:tiecompare} compares automatic scores on \emph{tie-good} vs.\ \emph{tie-bad} subsets.
COMET rises ${\approx}12$ points between groups, so metrics capture \emph{some} cultural signal.
However, \emph{tie-bad} sentences still average 66.2--66.9 COMET---only ${\approx}4$ below the global mean---understating failure severity.

\begin{table}[t]
  \centering
  \small
  \begin{tabular}{@{}lrrr@{}}
    \toprule
    & \textbf{\emph{tie-good}} & \textbf{\emph{tie-bad}} & \textbf{$\Delta$} \\
    \midrule
    GLM COMET & 78.95 & 66.22 & +12.72 \\
    DS COMET & 78.65 & 66.88 & +11.77 \\
    GLM chrF & 47.41 & 35.60 & +11.81 \\
    DS chrF & 46.16 & 38.23 & +7.92 \\
    \bottomrule
  \end{tabular}
  \caption{Mean automatic scores by human label subset.}
  \label{tab:tiecompare}
\end{table}

\paragraph{Metric--human correlation.}
Defining cultural adequacy score as the number of models passing (0--2), mean sentence COMET correlates weakly with human judgments (Spearman $\rho=0.29$).
On \emph{tie-good} vs.\ \emph{tie-bad} sentences alone, $\rho=0.45$.
On 116 \emph{g}/\emph{d} sentences with a unique human winner, COMET's segment-level winner matches humans in 54 cases (46.6\%); COMET margin (DeepSeek $-$ GLM) does not predict human preference ($\rho=-0.07$).
Among \emph{tie-bad} sentences, 50.0\% still assign higher COMET to DeepSeek-V4-Flash (mean DS COMET 66.88 vs.\ 70.87 globally).

\subsection{Reverse Misalignment Cases}

Aggregate statistics understate how often COMET \emph{rewards} the culturally inadequate output.
Table~\ref{tab:reverse} lists five curated cases from \texttt{reverse\_misalignment\_cases.csv} where human gold and COMET diverge most instructively---including \emph{inverse winners} (COMET favors the whitewashed hypothesis) and \emph{false-high} consensus failures (both models fail culturally yet COMET remains high).

\begin{table}[t]
  \centering
  \footnotesize
  \setlength{\tabcolsep}{3pt}
  \begin{tabular}{@{}clrrlp{4.8cm}@{}}
    \toprule
    \textbf{\#} & \textbf{Gold} & \textbf{GLM} & \textbf{DS} & \textbf{Kw.} & \textbf{Takeaway} \\
    \midrule
    1 & \emph{d} & 82.9 & 68.8 & 瓜娃子 & COMET prefers bland ``fool''; human accepts DS insult register \\
    2 & \emph{d} & 67.1 & 55.9 & 安逸 & COMET rewards ``comfortable''; human accepts DS pleasure reading \\
    3 & \emph{g} & 42.0 & 49.1 & 幺/九通 & Mahjong culture collapsed; COMET still prefers DS \\
    4 & \emph{tie-bad} & 77.8 & 81.7 & 咋个 & Both miss dialect address; COMET ${\approx}80$ \\
    5 & \emph{tie-bad} & 51.7 & 53.6 & 老子 & Both normalize self-reference; moderate COMET \\
    \bottomrule
  \end{tabular}
  \caption{Reverse misalignment cases: COMET ($\times 100$) vs.\ human cultural adequacy. GLM/DS = COMET scores.}
  \label{tab:reverse}
\end{table}

\noindent\textbf{Case 1} (\emph{d}, 瓜娃子): human judges accept only DeepSeek-V4-Flash for preserving insult force in ``treating me as 瓜娃子''; COMET favors GLM by +14.1.
\textbf{Case 2} (\emph{d}, 安逸): only DeepSeek conveys lived Sichuanese pleasure; COMET still rewards GLM's ``comfortable'' paraphrase (+11.2).
\textbf{Case 3} (\emph{g}, 幺/九通): mahjong term 九通 collapses to ``wine barrel'' in both outputs; COMET margin (+7.1) prefers DeepSeek despite human preference for GLM on 幺儿.
\textbf{Cases 4--5} (\emph{tie-bad}): both models miss 咋个伯父 (dialect address) and 老子 (emphatic self-reference) while COMET scores remain ${\approx}80$ and ${\approx}53$ respectively---fluent normalization masks pragmatic loss that humans label consensus failure.

\section{Discussion}

\subsection{Why COMET Creates an Illusion of Parity}
COMET optimizes semantic similarity to references.
Dialectal whitewashing often substitutes culturally rich terms with broadly synonymous standard terms occupying similar embedding space; COMET treats these as valid paraphrases rather than pragmatic failure (Cases 1--2, Table~\ref{tab:reverse}).
System-level aggregation further hides keyword-specific valleys (RQ2): 老子 and 安逸 score well below corpus means despite fluent outputs.

\subsection{Typology of Consensus Whitewashing}
Manual inspection of \emph{tie-bad} and contested \emph{g}/\emph{d} cases yields three recurring patterns:
\textbf{Pragmatic flattening}---老子 in imperatives (e.g., 给老子坐车) rendered as polite offers rather than emphatic demand (Case 5);
\textbf{literal standardization}---安逸 mapped to ``comfortable'' instead of embodied regional pleasure (Case 2);
\textbf{long-tail cultural collapse}---domain terms such as mahjong 九通 hallucinated as literal objects (``wine barrel''; Case 3).

\subsection{Implications for NLG Evaluation}
Variation-rich NLG should report stratified human cultural adequacy alongside reference-based metrics, especially for dialectal and personalised generation.
Shared \emph{tie-bad} failure (28\%) indicates normalization is often systemic rather than model-specific; evaluation protocols that only ask ``which translation is better'' cannot detect this.
We release our probe pipeline to support reproducible, variation-sensitive evaluation beyond corpus averages.

\section{Limitations}

Our study has several bounds:
(1)~\textbf{Scope}: one dialect, one direction, $N=200$, keyword-biased probe---not a random sample.
(2)~\textbf{Models}: two Chinese-strong LLMs only; no open- vs.\ closed-source generalization.
(3)~\textbf{Source noise}: WSC-Train ASR artifacts may conflate with model error in some \emph{tie-bad} labels.
(4)~\textbf{Reference bias}: a single culturally marked reference style affects COMET.
(5)~\textbf{Prompt}: one naive prompt may understate model capability.

\section{Conclusion}

System-level COMET parity between GLM-4-Plus and DeepSeek-V4-Flash on Sichuanese--English translation is misleading.
Human cultural adequacy reveals frequent consensus failure (28\% \emph{tie-bad}), significant paired pass-rate differences (McNemar $p=0.001$), and weak alignment with automatic metrics ($\rho=0.29$; 46.6\% winner agreement on contested items).
We argue for culturally informed, stratified NLG evaluation as a complement to standard reference-based scoring.

\section*{Ethics Statement}
Our probe uses publicly released speech transcriptions.
Human references and annotations were produced under internal guidelines.
Dialect examples include insult/register terms for linguistic analysis; we do not endorse pejorative language.

\section*{Reproducibility Statement}
Probe construction and evaluation are reproduced by \texttt{download\_sichuanese\_data.py}, \texttt{filter\_seeds.py} (seed 42; 45 keyword patterns; weighted sampling), \texttt{translate\_experiment.py} (prompts and API models as in \S\ref{sec:systems}), and \texttt{analyze\_scores.py} (sacreBLEU; COMET \texttt{wmt22-comet-da}; bootstrap $B=2000$).
Source: \texttt{ASLP-lab/WSC-Train}.
Outputs: \texttt{translation\_results\_200.csv}, \texttt{translation\_segment\_scores.csv}, \texttt{translation\_keyword\_scores.csv}.

\section*{Data Availability}
Annotations, model outputs, evaluation scores, and reproduction scripts are available at \url{https://github.com/yhanqi0611-create/sc-wsc-probe}.
Source transcriptions are from HuggingFace \texttt{ASLP-lab/WSC-Train}; the 200-sentence keyword-stratified probe is in \texttt{filtered\_seeds\_200.csv}.

\section*{Acknowledgments}
Omitted for review.

\bibliography{references}

\appendix

\section{Annotation Guidelines (Excerpt)}
\label{app:guidelines}

Annotators judge \textbf{cultural/pragmatic transmission}, not fluency alone.

\paragraph{安逸.}
\emph{Acceptable}: conveys spontaneous Sichuanese pleasure or satisfaction.
\emph{Reject}: generic ``comfortable'' with no regional pragmatic color.

\paragraph{老子.}
\emph{Acceptable}: emphatic first-person force in imperative/boast contexts.
\emph{Reject}: ``father'', ``Laozi'', or de-emphasized polite paraphrase.

\paragraph{瓜娃子.}
\emph{Acceptable}: preserves insult/register in context.
\emph{Reject}: neutral ``silly person'' that defuses aggression.

\end{document}
